\Article{25}
Besides the experiments just described, we carried out many others of the same kind, though we took much less care with the earlier ones. Still, it will be useful to bring their results together here. We leave out those obtained before the more refined instruments were set up, using cruder means and needles of widely differing dimensions, even though \textit{all} of them gave at least an approximation to the true value. The following values of $T$ were obtained in successive experiments.
\SourceLeaf{119}
\begin{center}
\begin{tabular}{c|l|l}
Number&Date, $1832$&$T$\\
I&May $21$&$1{,}7820$\\
II&May $24$&$1{,}7694$\\
III&June $4$&$1{,}7713$\\
IV&June $24$—$28$&$1{,}7625$\\
V&July $23$, $24$&$1{,}7826$\\
VI&July $25$, $26$&$1{,}7845$\\
VII&Sept. $9$&$1{,}7764$\\
VIII&Sept. $18$&$1{,}7821$\\
IX&Sept. $27$&$1{,}7965$\\
X&Oct. $15$&$1{,}7860$
\end{tabular}
\end{center}
Experiments V—IX were all carried out at the same place, whereas I—IV were carried out at other places. Experiment X is really a mixed case, since the deflections were observed at the usual place and the oscillations elsewhere. Almost equal care was taken with experiments VII and VIII, a little less with IV, V, VI, X, and much less with I—III. Experiments I—VIII used different needles of nearly the same weight and length (the weight was between $400$ and $440$ grams), whereas experiment X used a needle weighing $1062$ grams and measuring $485$ millimeters in length. Experiment IX was undertaken solely to see what degree of precision can be attained with a rather small needle. The needle used weighed only $58$ grams, but otherwise we took no less care than in experiments VII and VIII. There is no doubt that the precision of the observations will increase noticeably if still heavier needles are used, for example, needles weighing up to $2000$ or $3000$ grams.

